\section{LlamaIndex 源码分析：Query 流程}

\subsection{检索过程}

\begin{enumerate}
    \item 对用户 Query 计算 Embedding Vector（同样 1536 维）
    \item 计算 Query Embedding 与所有 Node Embedding 的相似度
    \item 默认取 \textbf{Top-2} 最相关的 Node
\end{enumerate}

相似度计算支持多种度量：欧式距离、点积、余弦相似度等。

\subsection{Prompt Template}

\begin{importantbox}{QA Template（源码中提取）}
核心模板：
\begin{quote}
\textit{Context information is below. \{context\} Given the context information and not prior knowledge, answer the query. Query: \{query\}}
\end{quote}
关键约束：\textbf{仅基于 Context Information 回答，不使用模型已有知识。}
\end{importantbox}

此外还有 Refine Template（基于新 Context 优化已有答案）和 Summarize Template 等。

\subsection{消息构建与 API 调用}

最终构建的消息包含：
\begin{enumerate}
    \item \textbf{System Message}：角色设定——"你是一个值得信赖的 QA 专家，永远基于 Context Information 回答，不使用已有知识"
    \item \textbf{User Message}：包含 Context Information（Top-K 相关文档）+ Query
\end{enumerate}

默认 Response Mode 为 Compact（紧凑型）。

\subsection{本章小结}

Query 流程：Query Embedding $\rightarrow$ 相似度匹配 $\rightarrow$ Top-K Node $\rightarrow$ Prompt Template 组装 $\rightarrow$ LLM 生成回答。整个过程只有约五行框架代码，但内部源码非常精细。

